\clearpage
\setcounter{chapter}{9}
\chaptermark{吸收、色散、散射}
\pdfbookmark[0]{第九章 吸收、色散、散射}{chapter9}
\problem{9-1}{
有一介质其吸收系数 $\alpha$ 为 $0.32\,\mathrm{cm}^{-1}$，透射光强衰减为入射光强的 $10\%$、$20\%$、$50\%$ 和 $80\%$ 时，相应的介质厚度各为多少？
}
\problem{9-2}{
一玻璃管长 $3.50\,\mathrm m$，内存有标准大气压下的某种气体，其吸收系数为 $0.1650\,\mathrm m^{-1}$。
\begin{enumerate}
\item 若仅考量这气体的吸收，求出透射光强的百分比。
\item 若再考量管口玻璃表面的反射，求出透射光强的百分比。设此玻璃的光强反射率为 $4\%$，并忽略多次反射和干涉。
\end{enumerate}
}
\problem{9-3}{
某种无色透明玻璃的吸收系数为 $0.10\,\mathrm m^{-1}$，用以制成 $5\,\mathrm{mm}$ 厚的玻璃窗。
\begin{enumerate}
\item 若仅考量这玻璃的吸收，求出透射光强的百分比。
\item 若再考量玻璃表面的反射，求出透射光强的百分比。设此玻璃的光强反射率为 $4\%$，并忽略多次反射和干涉。
\end{enumerate}
}
\problem{9-4}{
人眼能觉察的光强是太阳到达地面光强的 $1/10^{18}$，试问人在海底多少深度还能看见亮光？设海水吸收系数为 $1.0\,\mathrm m^{-1}$。
}
\problem{9-5}{
某玻璃对氦氖激光 $633\,\mathrm{nm}$ 的复折射率为 $\widetilde n=1.5+5\times10^{-8}\mathrm i$，求出该玻璃的吸收系数，以及这激光束在玻璃中的光速。
}
\problem{9-6}{
1908 年米氏以导体小球为模型建立了微粒对电磁波的散射理论。他设这小球的复折射率为 $\widetilde n=0.57+2.45\mathrm i$，针对真空中波长为 $550\,\mathrm{nm}$ 的光波，这导体材料的吸收系数为多少？若以光强衰减为入射光强的 $1/\mathrm e^2$ 作为穿透深度的定义，试求光在这导体小球中的透射深度 $d$ 为多少？
}
\problem{9-7}{
水银灯光含有两条显著的谱线，一条蓝色 $\lambda=435.8\,\mathrm{nm}$，另一条绿色 $\lambda=546.1\,\mathrm{nm}$。某一光学玻璃对这两波长的折射率分别为 $1.6525$ 和 $1.6245$。
\begin{enumerate}
\item 试根据以上数据确定出柯西公式的两个常数 $A$ 和 $B$。
\item 推算出这种玻璃对钠黄光 $\lambda=589.3\,\mathrm{nm}$ 的折射率。
\item 进一步导出这种玻璃在钠黄光附近的色散率 $\mathrm dn/\mathrm d\lambda$。
\item 钠黄光为双线结构，含 $\lambda_1=589.0\,\mathrm{nm}$ 和 $\lambda_2=589.6\,\mathrm{nm}$。试估算这两种波长的光在该玻璃中的相速之差 $\Delta v$ 与平均相速 $\bar v$ 之比值（数量级）；同时算出双线所形成的波包的群速 $v_{\mathrm g}$ 及其与 $\bar v$ 之比值（数量级）。
\end{enumerate}
}
\problem{9-8}{
依据正文图 9.3（b）中显示的水晶的正常色散曲线 $n(\lambda)$，
\begin{enumerate}
\item 较精确地定出水晶对紫光 $\lambda=400\,\mathrm{nm}$ 的折射率和色散率 $\mathrm dn/\mathrm d\lambda$。
\item 若一波包其中心波长为 $400\,\mathrm{nm}$，它传播于水晶中的群速 $v_{\mathrm g}$ 为多少？
\end{enumerate}
}
\problem{9-9}{
用折射法测得二硫化碳的折射率为：
\[
\lambda=527\,\mathrm{nm},\ n=1.642;\qquad
\lambda=589\,\mathrm{nm},\ n=1.629;
\]
\[
\lambda=656\,\mathrm{nm},\ n=1.620.
\]
试求波长为 $589\,\mathrm{nm}$ 附近的光在二硫化碳中的相速 $v_{\mathrm p}$ 和群速 $v_{\mathrm g}$。
}
\problem{9-10}{
波长为 $0.67\,\mathrm{nm}$ 的 X 射线由真空入射到某种玻璃时，不发生全反射的最小掠射角为 $0.1^\circ$。试确定该玻璃对这 X 射线的折射率。
}
\problem{9-11}{
一棱镜的顶角为 $50^\circ$，其所用的玻璃材料的色散可由二常数的柯西公式给予描述，且 $A=1.54$，$B=4.653\times10^3\,\mathrm{nm^2}$。
\begin{enumerate}
\item 求此棱镜对波长为 $550\,\mathrm{nm}$ 光束的最小偏向角 $\delta_{\mathrm m}$。
\item 求对此波长且调节到 $\delta_{\mathrm m}$ 条件下该棱镜的角色散 $\mathrm d\delta_{\mathrm m}/\mathrm d\lambda$。
\item 设该棱镜主截面的底边长 $B_0$ 为 $6\,\mathrm{cm}$，求出该棱镜对 $550\,\mathrm{nm}$ 光的色分辨本领 $\lambda/\Delta\lambda_{\mathrm m}$，以及相应的可分辨的最小波长间隔 $\delta\lambda_{\mathrm m}$。
\end{enumerate}
}
\problem{9-12}{
自由电子气对电磁波的色散关系为
\[
n^2(\omega)=1-\frac{\omega_{\mathrm p}^2}{\omega^2},
\]
其中 $\omega_{\mathrm p}$ 为一特征频率，可视为一常数。
\begin{enumerate}
\item 试导出其色散关系的又一形式
\[
\omega^2(k)=c^2k^2+\omega_{\mathrm p}^2.
\]
\item 进而证明电磁波在自由电子气中的相速 $v$ 与群速 $v_{\mathrm g}$ 之乘积为一常数，
\[
v\cdot v_{\mathrm g}=c^2.\qquad (c\text{ 为真空中光速})
\]
\end{enumerate}
}
\problem{9-13}{
正文第 5 节以谱密度函数为高斯线型的波包作为对象，并计及二阶色散效应 $\ddot\omega$，而逐一导出波包中心速度 $v_0$、波包展宽 $\Delta x(t)$ 和波包前沿速度即讯号速度 $v_{\mathrm i}$。建议试以谱密度函数为方盒线型的波包作为对象，依次导出 $v_0$、$\Delta x(t)$ 和 $v_{\mathrm i}$ 公式。
}
\problem{9-14}{
用氩离子激光其波长为 $488.0\,\mathrm{nm}$ 照射金刚石，发现其散射光的波长中除 $488.0\,\mathrm{nm}$ 谱线外，还有 $\lambda'=458.21\,\mathrm{nm}$ 和 $\lambda'=521.92\,\mathrm{nm}$ 两条拉曼谱线。问，金刚石的拉曼特征频移为多少（用波数 $\mathrm{cm}^{-1}$ 表示）？
}
\problem{9-15}{
苯 $(\mathrm C_6\mathrm H_6)$ 的拉曼光谱中较强的谱线与入射光的波数差分别有 $607\,\mathrm{cm}^{-1}$、$992\,\mathrm{cm}^{-1}$、$1178\,\mathrm{cm}^{-1}$、$1586\,\mathrm{cm}^{-1}$、$3047\,\mathrm{cm}^{-1}$ 和 $3062\,\mathrm{cm}^{-1}$。现以氩离子激光其波长为 $488\,\mathrm{nm}$ 入射，试计算靠近 $488\,\mathrm{nm}$ 谱线两侧前三个斯托克斯谱线和反斯托克斯谱线的波长。
}
\problem{9-16}{
一个长 $30\,\mathrm{cm}$ 的管中充有含烟的气体，它能透过 $60\%$ 的光强，而将烟尘完全去除后，则可透过 $92\%$ 的光强。设烟尘对光只有散射而无吸收，试求出气体的吸收系数 $\alpha_a$ 和烟尘的散射系数 $\alpha_s$。
}
\problem{9-17}{
试在坐标纸上绘制出一条反映瑞利散射强度与波长关系的曲线，建议横坐标上波长值分别取 $300\,\mathrm{nm}$、$350\,\mathrm{nm}$、$400\,\mathrm{nm}$、$450\,\mathrm{nm}$、$550\,\mathrm{nm}$、$650\,\mathrm{nm}$、$750\,\mathrm{nm}$；纵坐标上选取 $300\,\mathrm{nm}$ 时的瑞利散射强度（相对值）为 $1$。
}
